\chapter{Complex Numbers}

\begin{remark}
In engineering, the imaginary unit is written $j$ (rather than $i$) to avoid
confusion with electric current.
\end{remark}

% -------------------------------------------------------
\section{Definition}
% -------------------------------------------------------

\begin{definition}[Complex Number]
A \textbf{complex number} is an expression of the form
\[
  z = x + jy,
\]
where $x = \operatorname{Re}(z)$ is the \textbf{real part}, $y = \operatorname{Im}(z)$ is the
\textbf{imaginary part}, and $j = \sqrt{-1}$.
\end{definition}

% -------------------------------------------------------
\section{Argand Diagram}
% -------------------------------------------------------

\begin{definition}[Argand Diagram]
The \textbf{Argand diagram} is a Cartesian plane used to represent complex numbers
geometrically:
\begin{itemize}
  \item the horizontal axis represents the real part,
  \item the vertical axis represents the imaginary part.
\end{itemize}
The complex number $z = x + jy$ corresponds to the point $(x, y)$.
\end{definition}

% -------------------------------------------------------
\section{Modulus and Argument}
% -------------------------------------------------------

\begin{definition}[Modulus]
The \textbf{modulus} (or absolute value) of $z = x + jy$ is the distance from the
origin to the point $(x,y)$:
\[
  |z| = r = \sqrt{x^2 + y^2}.
\]
\end{definition}

\begin{definition}[Argument]
The \textbf{argument} of $z$ is the angle $\theta$ that the line $Oz$ makes with the
positive real axis:
\[
  \theta = \arg(z) = \tan^{-1}\!\left(\frac{y}{x}\right).
\]
The \textbf{principal argument} satisfies $-\pi < \theta \leq \pi$.
\end{definition}

% -------------------------------------------------------
\section{Polar and Exponential Forms}
% -------------------------------------------------------

\begin{definition}[Polar Form]
Using $x = r\cos\theta$ and $y = r\sin\theta$, any complex number can be written as
\[
  z = r(\cos\theta + j\sin\theta).
\]
\end{definition}

\begin{definition}[Exponential Form]
By Euler's formula (Section~\ref{sec:euler}), the polar form condenses to
\[
  z = re^{j\theta}.
\]
\end{definition}

\begin{note}
The conversion between forms is:
\[
  x = r\cos\theta, \qquad y = r\sin\theta, \qquad
  r = \sqrt{x^2+y^2}, \qquad \theta = \tan^{-1}\!\!\left(\frac{y}{x}\right).
\]
\end{note}

% -------------------------------------------------------
\section{Multiplication and Division}
% -------------------------------------------------------

\begin{theorem}[Multiplication in Polar Form]
If $z_1 = r_1 e^{j\theta_1}$ and $z_2 = r_2 e^{j\theta_2}$, then
\[
  z_1 z_2 = r_1 r_2\,\bigl[\cos(\theta_1+\theta_2) + j\sin(\theta_1+\theta_2)\bigr].
\]
Moduli multiply; arguments add.
\end{theorem}

\begin{theorem}[Division in Polar Form]
\[
  \frac{z_1}{z_2} = \frac{r_1}{r_2}\,\bigl[\cos(\theta_1-\theta_2) + j\sin(\theta_1-\theta_2)\bigr].
\]
Moduli divide; arguments subtract.
\end{theorem}

% -------------------------------------------------------
\section{De Moivre's Theorem}
% -------------------------------------------------------

\begin{theorem}[De Moivre's Theorem]
For any integer $n$:
\[
  z^n = r^n(\cos n\theta + j\sin n\theta).
\]
\end{theorem}

\begin{theorem}[$n$-th Roots]
The $n$ distinct $n$-th roots of $z = re^{j\theta}$ are
\[
  z_k = r^{1/n}\left[\cos\!\left(\frac{\theta + 2k\pi}{n}\right)
                    + j\sin\!\left(\frac{\theta + 2k\pi}{n}\right)\right],
  \qquad k = 0, 1, \ldots, n-1.
\]
The roots are equally spaced around a circle of radius $r^{1/n}$ in the Argand diagram.
\end{theorem}

% -------------------------------------------------------
\section{Euler's Formula and Key Identities}
\label{sec:euler}
% -------------------------------------------------------

\begin{theorem}[Euler's Formula]
\[
  e^{j\theta} = \cos\theta + j\sin\theta.
\]
\end{theorem}

From Euler's formula, adding and subtracting $e^{j\theta}$ and $e^{-j\theta}$ yields:

\begin{corollary}
\[
  \cos\theta = \frac{e^{j\theta} + e^{-j\theta}}{2}, \qquad
  \sin\theta = \frac{e^{j\theta} - e^{-j\theta}}{2j}.
\]
\end{corollary}

\begin{example}[Euler's Identity]
Setting $\theta = \pi$:
\[
  e^{j\pi} + 1 = 0.
\]
\end{example}
